\documentclass[11pt]{article}
\usepackage[left=2cm, right=2cm, top=2cm, bottom = 2cm]{geometry}

\usepackage{amsmath, enumitem, tcolorbox, hyperref}

% Fill in your name below!!
\title{TR Title}
\author{Your Names Here}
\date{\today}

\begin{document}

\maketitle


\begin{abstract}
    Include your abstract here. Abstracts should typically be between 100-200 words. One option for structuring an abstract is to begin with a sentence on background, then proceed to the objective of the report, your unique contributions to the problem (aka what you did), and briefly summarize results / findings.
\end{abstract}

\section*{Background and Objective}

Throughout the entire report, write in \textit{complete sentences}, and make sure that each paragraph has a \textit{topic sentence}. \textbf{BE CONCISE}. A longer report does not directly correspond to a \textit{better} report. Don't ramble. Say what you need to say (clearly), and move on. The same goes for the following sections.

\noindent Things that should be included in this section:

\begin{itemize}
    \item What is the motivating context?
    \item What question(s) are you trying to answer?
    \item If there are \textit{relevant} summary measures or visualizations, include them here (as tables or figures), and make sure they are labeled \textit{and} have captions.
\end{itemize}



\section*{Methodology}

\begin{itemize}
    \item What sort of model did you use to address your questions?
    \item What data decisions (if any) did you make that might impact your results, or impacted your modeling decisions?
    \item With \textit{rare} exceptions, this section should not contain visualizations or tables.
    \item If you used code / software for your analysis, what did you use? What packages / versions? Reference a code file that you submit along with your report so that someone could replicate your analysis.
    \item If there are details that would be relevant to replicating your analysis but may \textit{detract} from the main ``point" of your report, considering including them in an appendix and referencing sections of it here.
    \item Justify your modeling choices, either using context and/or statistical theory. For example, if your outcome is binary, a binomial likelihood would be theoretically justified. If you're interested in \textit{change} in a quantitative variable over time, you could justify (1) calculating a difference, and modeling that difference with a continuous random variable, OR (2) modeling the latter observation as a function of the first, so that your results are interpreted for observations \textit{conditional on} a given initial value.
\end{itemize}

\noindent Include any model statements in \textit{formal notation} as equations. If you plan on referencing them in your report, label them (using \texttt{align} rather than \texttt{align*}). Here's an example of how you might write up a model statement appropriately. \\

%% Example start
\noindent Let $y_i$ denote the number of deaths out of $n_i$ total individual in county $i = 1, \dots, 400$. We model the probability of death in each county $\theta_i$ as a linear function of observed average county income $x_i$

\begin{align}
\begin{split} % split makes sure that the equation is only labeled ONCE. if you want each line of the equation to have a different number, remove the split environment and use only align!
    y_i \mid \theta_i, n_i & \sim \text{Binomial}(n_{i}, \theta_i) \\
    \theta_i & = \beta_0 + \beta_1 x_i \\
    \beta_0 & \propto 1 \\
    \beta_1 & \sim N(0, 1)
\end{split}
\end{align}
where $\beta_0$ is given an improper priors, and $\beta_1$ is given a standard normal prior.
%% Example end

\noindent General rules of thumb:
\begin{itemize}
    \item Vectors should be bolded. Use \texttt{textbf} or \texttt{boldsymbol} if you're bolding a letter or symbol, respectively.
    \item Observations should be indexed according to how they vary. In the above example, we have 400 observations (1 per county), so we chose the index $i$ to represent county. If we additionally had observations across \textit{time} and a covariate that only varied by time, for example, we would need an additional index for time $t$.
\end{itemize}

\section*{Analysis and Results}

\begin{itemize}
    \item \textit{Report} your results. Make a note of summary measures, interpret credible intervals, convergence diagnostics, and posterior predictive checks, as relevant for the report.
    \item Include figures and tables of relevant summaries. This might include trace plots, histograms / boxplots / density plots of samples from posterior distributions, etc.
    \item Do not discuss the \textit{impact} of your results! That comes in the discussion section.
\end{itemize}

\section*{Discussion}

\begin{itemize}
    \item Answer the question you were interested in all along, based on your results.
    \item Discuss any limitations of your analysis.
    \item What would/could be included in a \textit{future} analysis or futher study?
\end{itemize}

%% References

%% Appendix



\end{document}